\documentclass[10pt,a4paper]{amsart}
\usepackage{amsmath,amssymb,amsthm,graphicx}
% Retained minimal preamble for the Stacks Project chapter more-morphisms.tex (More on Morphisms).
% Verbatim excerpt of the chapter preamble: Xy-pic setup (native preamble lines 14-18) and the native
% theorem declarations with the mathematical macros (native preamble lines 154-217). Copied mechanically;
% nothing added, removed or redefined. The wrapper class and the amsmath/amssymb/amsthm/graphicx base
% packages are supplied by the assembler, so the class/externaldocument part of the native preamble is omitted.
% For commutative diagrams we use Xy-pic
\usepackage[all]{xy}

% We use 2cell for 2-commutative diagrams.
\xyoption{2cell}
% Theorem environments.
%
\theoremstyle{plain}
\newtheorem{theorem}[subsection]{Theorem}
\newtheorem{proposition}[subsection]{Proposition}
\newtheorem{lemma}[subsection]{Lemma}

\theoremstyle{definition}
\newtheorem{definition}[subsection]{Definition}
\newtheorem{example}[subsection]{Example}
\newtheorem{exercise}[subsection]{Exercise}
\newtheorem{situation}[subsection]{Situation}

\theoremstyle{remark}
\newtheorem{remark}[subsection]{Remark}
\newtheorem{remarks}[subsection]{Remarks}

\numberwithin{equation}{subsection}

% Macros
%
\def\lim{\mathop{\mathrm{lim}}\nolimits}
\def\colim{\mathop{\mathrm{colim}}\nolimits}
\def\Spec{\mathop{\mathrm{Spec}}}
\def\Hom{\mathop{\mathrm{Hom}}\nolimits}
\def\Ext{\mathop{\mathrm{Ext}}\nolimits}
\def\SheafHom{\mathop{\mathcal{H}\!\mathit{om}}\nolimits}
\def\SheafExt{\mathop{\mathcal{E}\!\mathit{xt}}\nolimits}
\def\Sch{\mathit{Sch}}
\def\Mor{\mathop{\mathrm{Mor}}\nolimits}
\def\Ob{\mathop{\mathrm{Ob}}\nolimits}
\def\Sh{\mathop{\mathit{Sh}}\nolimits}
\def\NL{\mathop{N\!L}\nolimits}
\def\CH{\mathop{\mathrm{CH}}\nolimits}
\def\proetale{{pro\text{-}\acute{e}tale}}
\def\etale{{\acute{e}tale}}
\def\QCoh{\mathit{QCoh}}
\def\Ker{\mathop{\mathrm{Ker}}}
\def\Im{\mathop{\mathrm{Im}}}
\def\Coker{\mathop{\mathrm{Coker}}}
\def\Coim{\mathop{\mathrm{Coim}}}

% Boxtimes
%
\DeclareMathSymbol{\boxtimes}{\mathbin}{AMSa}{"02}

%
% Macros for moduli stacks/spaces
%
\def\QCohstack{\mathcal{QC}\!\mathit{oh}}
\def\Cohstack{\mathcal{C}\!\mathit{oh}}
\def\Spacesstack{\mathcal{S}\!\mathit{paces}}
\def\Quotfunctor{\mathrm{Quot}}
\def\Hilbfunctor{\mathrm{Hilb}}
\def\Curvesstack{\mathcal{C}\!\mathit{urves}}
\def\Polarizedstack{\mathcal{P}\!\mathit{olarized}}
\def\Complexesstack{\mathcal{C}\!\mathit{omplexes}}
% \Pic is the operator that assigns to X its picard group, usage \Pic(X)
% \Picardstack_{X/B} denotes the Picard stack of X over B
% \Picardfunctor_{X/B} denotes the Picard functor of X over B
\def\Pic{\mathop{\mathrm{Pic}}\nolimits}
\def\Picardstack{\mathcal{P}\!\mathit{ic}}
\def\Picardfunctor{\mathrm{Pic}}
\def\Deformationcategory{\mathcal{D}\!\mathit{ef}}
\newenvironment{reference}{}{}
\newenvironment{slogan}{}{}
\newenvironment{history}{}{}
\usepackage{cleveref}
\crefname{originallemma}{original lemma}{original lemmas}
\Crefname{originallemma}{Original lemma}{Original lemmas}
\makeatletter
\newlabel{lemma-quasi-finite-separated-quasi-affine}{{02LR}{source}{}{sourcepointer-1}{}}
\newlabel{lemma-quasi-finite-separated-quasi-affine@cref}{{[originallemma][0][]02LR}{[1][0][]source}{}{}{}}
\makeatother
\crefname{originallemma}{original lemma}{original lemmas}
\Crefname{originallemma}{Original lemma}{Original lemmas}
\makeatletter
\newlabel{lemma-quasi-projective}{{0B42}{source}{}{sourcepointer-2}{}}
\newlabel{lemma-quasi-projective@cref}{{[originallemma][0][]0B42}{[1][0][]source}{}{}{}}
\makeatother
\crefname{originallemma}{original lemma}{original lemmas}
\Crefname{originallemma}{Original lemma}{Original lemmas}
\makeatletter
\newlabel{morphisms-lemma-base-change-quasi-projective}{{0B3G}{source}{}{sourcepointer-3}{}}
\newlabel{morphisms-lemma-base-change-quasi-projective@cref}{{[originallemma][0][]0B3G}{[1][0][]source}{}{}{}}
\makeatother
\crefname{originallemma}{original lemma}{original lemmas}
\Crefname{originallemma}{Original lemma}{Original lemmas}
\makeatletter
\newlabel{morphisms-lemma-immersion-locally-quasi-finite}{{01TN}{source}{}{sourcepointer-4}{}}
\newlabel{morphisms-lemma-immersion-locally-quasi-finite@cref}{{[originallemma][0][]01TN}{[1][0][]source}{}{}{}}
\makeatother
\crefname{originallemma}{original lemma}{original lemmas}
\Crefname{originallemma}{Original lemma}{Original lemmas}
\makeatletter
\newlabel{morphisms-lemma-projective-quasi-projective}{{07RL}{source}{}{sourcepointer-5}{}}
\newlabel{morphisms-lemma-projective-quasi-projective@cref}{{[originallemma][0][]07RL}{[1][0][]source}{}{}{}}
\makeatother
\crefname{originallemma}{original lemma}{original lemmas}
\Crefname{originallemma}{Original lemma}{Original lemmas}
\makeatletter
\newlabel{morphisms-lemma-quasi-affine-finite-type-quasi-projective}{{0B3H}{source}{}{sourcepointer-6}{}}
\newlabel{morphisms-lemma-quasi-affine-finite-type-quasi-projective@cref}{{[originallemma][0][]0B3H}{[1][0][]source}{}{}{}}
\makeatother
\title{The quasi-projective schemes over a base form a category closed under base change, composition, fibre products, sums, and standard sources}
\author{The Stacks Project Authors}
\date{stacks-project:more-morphisms.tex:local-snapshot-20261011; GNU FDL 1.2 or later}
\begin{document}
\maketitle
\noindent\textit{Editorial note.} Statement span is lines 13891-13914 of sources/stacks-more-morphisms/pinned-original-source.tex; proof span is lines 13916-13954, the complete author proof (37 body lines).\par\medskip
\noindent\textit{Editorial note.} Navigation stubs: 6 cited-but-not-included labels are supplied as non-load-bearing pointers with their pinned tags where the snapshot tag index has them.\par\medskip
\noindent\textit{Editorial note.} extra\_packages is empty on purpose: the retained preamble already loads xy with options and a second declaration clashes in the assembler wrapper.\par\medskip
\noindent\textit{Editorial note.} Text is retained verbatim, including the author's own typos and placeholders.\par\medskip
\noindent\textit{Editorial note.} Companion to the already-collected projective-category card L16075; the placeholder list item is part of the pinned statement and was not silently dropped.\par\medskip
\noindent\textit{Allowed standard prerequisites (editorial).}\begin{itemize}
\item Explicit prerequisite boundary. The quasi-projective five-way equivalence (lemma-quasi-projective, collected as L16072); the Morphisms-chapter base change and permanence lemmas for quasi-projective morphisms; the chapter's quasi-finite-separated-quasi-affine theorem (Zariski Main Type corollary).
\end{itemize}\medskip
\ifdefined\hypertarget\hypertarget{sourcepointer-1}{}\fi
\noindent\textit{Source reference.} Original lemma 02LR: sources/stacks-more-morphisms/pinned-original-source.tex, source lines 12426--12429. Cited by the retained statement or proof but not part of the collected spans; supplied only as a navigation pointer. The argument does not depend on its content.\par\medskip
\ifdefined\hypertarget\hypertarget{sourcepointer-2}{}\fi
\noindent\textit{Source reference.} Original lemma 0B42: sources/stacks-more-morphisms/pinned-original-source.tex, source lines 13840--13853. Cited by the retained statement or proof but not part of the collected spans; supplied only as a navigation pointer. The argument does not depend on its content.\par\medskip
\ifdefined\hypertarget\hypertarget{sourcepointer-3}{}\fi
\noindent\textit{Source reference.} Original lemma 0B3G: sources/stacks-more-morphisms/pinned-original-source.tex, source lines 13954--13954. Cited by the retained statement or proof but not part of the collected spans; supplied only as a navigation pointer. The argument does not depend on its content.\par\medskip
\ifdefined\hypertarget\hypertarget{sourcepointer-4}{}\fi
\noindent\textit{Source reference.} Original lemma 01TN: sources/stacks-more-morphisms/pinned-original-source.tex, source lines 13954--13954. Cited by the retained statement or proof but not part of the collected spans; supplied only as a navigation pointer. The argument does not depend on its content.\par\medskip
\ifdefined\hypertarget\hypertarget{sourcepointer-5}{}\fi
\noindent\textit{Source reference.} Original lemma 07RL: sources/stacks-more-morphisms/pinned-original-source.tex, source lines 13954--13954. Cited by the retained statement or proof but not part of the collected spans; supplied only as a navigation pointer. The argument does not depend on its content.\par\medskip
\ifdefined\hypertarget\hypertarget{sourcepointer-6}{}\fi
\noindent\textit{Source reference.} Original lemma 0B3H: sources/stacks-more-morphisms/pinned-original-source.tex, source lines 13954--13954. Cited by the retained statement or proof but not part of the collected spans; supplied only as a navigation pointer. The argument does not depend on its content.\par\medskip
\begin{lemma}
\label{lemma-category-quasi-projective}
Let $S$ be a scheme which has an ample invertible sheaf.
Let $\text{QP}_S$ be the full subcategory of the
category of schemes over $S$ satisfying the equivalent
conditions of Lemma \ref{lemma-quasi-projective}.
\begin{enumerate}
\item if $S' \to S$ is a morphism of schemes and $S'$ has
an ample invertible sheaf, then base change determines
a functor $\text{QP}_S \to \text{QP}_{S'}$,
\item if $X \in \text{QP}_S$ and $Y \in \text{QP}_X$, then $Y \in \text{QP}_S$,
\item the category $\text{QP}_S$ is closed under fibre products,
\item the category $\text{QP}_S$ is closed under
finite disjoint unions,
\item if $X \to S$ is projective, then $X \in \text{QP}_S$,
\item if $X \to S$ is quasi-affine of finite type, then
$X$ is in $\text{QP}_S$,
\item if $X \to S$ is quasi-finite and separated, then
$X \in \text{QP}_S$,
\item if $X \to S$ is a quasi-compact immersion, then
$X \in \text{QP}_S$,
\item add more here.
\end{enumerate}
\end{lemma}

\begin{proof}
Part (1) follows from Morphisms, Lemma
\ref{morphisms-lemma-base-change-quasi-projective}.

\medskip\noindent
Part (2) follows from the fourth characterization of
Lemma \ref{lemma-quasi-projective}.

\medskip\noindent
If $X \to S$ and $Y \to S$ are quasi-projective, then
$X \times_S Y \to Y$ is quasi-projective by
Morphisms, Lemma \ref{morphisms-lemma-base-change-quasi-projective}.
Hence (3) follows from (2).

\medskip\noindent
If $X = Y \amalg Z$ is a disjoint union of schemes
and $\mathcal{L}$ is an invertible $\mathcal{O}_X$-module
such that $\mathcal{L}|_Y$ and $\mathcal{L}|_Z$ are ample, then
$\mathcal{L}$ is ample (details omitted). Thus
part (4) follows from the fourth characterization of
Lemma \ref{lemma-quasi-projective}.

\medskip\noindent
Part (5) follows from
Morphisms, Lemma \ref{morphisms-lemma-projective-quasi-projective}.

\medskip\noindent
Part (6) follows from
Morphisms, Lemma
\ref{morphisms-lemma-quasi-affine-finite-type-quasi-projective}.

\medskip\noindent
Part (7) follows from part (6) and
Lemma \ref{lemma-quasi-finite-separated-quasi-affine}.

\medskip\noindent
Part (8) follows from part (7) and
Morphisms, Lemma \ref{morphisms-lemma-immersion-locally-quasi-finite}.
\end{proof}

\end{document}
